\documentclass[11pt,a4paper]{article}
\usepackage[T1]{fontenc}
\usepackage{lmodern,amsmath,amssymb,amsthm,mathtools,microtype}
\usepackage[margin=27mm]{geometry}
\usepackage[colorlinks=true,linkcolor=blue,citecolor=blue,urlcolor=blue]{hyperref}
\newtheorem{theorem}{Theorem}[section]
\newtheorem{lemma}[theorem]{Lemma}
\newtheorem{proposition}[theorem]{Proposition}
\newtheorem{corollary}[theorem]{Corollary}
\newcommand{\Q}{\mathbb Q}
\DeclareMathOperator{\Tr}{Tr}
\title{Rational Normalization of Local Valuation Tests}
\author{}\date{10 November 2026}
\begin{document}\maketitle
\begin{abstract}
We give rational functions which suppress local conditions at a prime
where an auxiliary coefficient has odd valuation, without requiring any
bound on the poles of the original functions. At a prime where that
coefficient is a unit, positive valuation of the normalized function is
equivalent to positive valuation of the original one. A second result
uses two balanced contact functions with disjoint contact and pole sets:
after a suitable change of scale, at least one is a unit at each prime
of odd valuation of the scalar coefficient. Together these results
supply local extraction statements for products of two arithmetic
coefficients. The statements concern local conditions and do not by
themselves establish a smaller defining polynomial for the integers.
\end{abstract}

\section{A rational function with controlled valuations}

Let $L$ be a finite extension of $\Q_p$ unramified over $\Q_p$.
Normalize its valuation by $v(p)=1$, so that $v(L^\times)=\mathbb Z$,
and put $v(0)=+\infty$. For $z\in L^\times$, define
\[
                   \Psi_z(D)=\frac{zD^2}{1+zD^2}
\]
whenever the denominator is nonzero. All assertions involving this
rational function include that domain condition.

\begin{lemma}\label{lem:normalization}
The following statements hold.
\begin{enumerate}
\item If $v(z)$ is odd, then the denominator of $\Psi_z(D)$ is
nonzero for every $D\in L$, and
\[
 v(\Psi_z(D))=\max\{0,v(z)+2v(D)\}\ge0
 \qquad(D\ne0).
\]
For $D=0$ the value of the rational function is zero.
\item If $v(z)=0$, then
\[
        v(\Psi_z(D))>0\quad\Longleftrightarrow\quad v(D)>0.
\]
When $0<v(D)<+\infty$, one has
$v(\Psi_z(D))=2v(D)$.
\end{enumerate}
\end{lemma}
\begin{proof}
For $D\ne0$ write $e=v(z)+2v(D)$. If $v(z)$ is odd, then
$e$ is odd and in particular is not zero. Thus
$v(1+zD^2)=\min(0,e)$, and subtraction gives the first formula.
For $D=0$ the assertions follow directly.

Suppose now that $v(z)=0$. If $v(D)<0$, the same computation
gives valuation zero. If $v(D)>0$, the denominator is a unit and
the valuation is $2v(D)>0$. Finally, if $v(D)=0$, the numerator
is a unit and the denominator has nonnegative valuation, so the
quotient has nonpositive valuation. This proves the equivalence,
including the zero value of $D$.
\end{proof}

In particular, if $v(z)=\nu>0$ is odd and $K$ is a nonnegative
integer, then
\begin{equation}\label{eq:inactive}
             v\bigl(z^K\Psi_z(D)\bigr)\ge K\nu
\end{equation}
for every $D$. The assertion is uniform in $D$: a pole of
arbitrarily large order does not weaken the lower bound. In contrast,
multiplication by $z^K$ alone cannot give such a bound, since
$v(z^KD)=K\nu+v(D)$ is unbounded below as $D$ varies.

\section{Two collections of local conditions}

Let $F/\Q$ be a number field and let $p$ be unramified in $F$.
For a place $\mathfrak p$ above $p$, normalize
$v_{\mathfrak p}(p)=1$. Let $z,w\in\Q^\times$, and let
$C_1,\ldots,C_a,S_1,\ldots,S_b$ be elements of $F$.
Fix positive integers $m_i,n_j$, and put
\[
        A_i=2m_i+1,\qquad B_j=2n_j+1.
\]
Choose an integer
\[
 K\ge1+\max\{A_1,\ldots,A_a,B_1,\ldots,B_b\}.
\]
Empty collections can simply be omitted. Define
\[
 R=zw,\qquad
 E_i=z^K\Psi_z(C_i),\qquad
 G_j=w^K\Psi_w(S_j),
\]
and
\begin{equation}\label{eq:product}
 W=\prod_{i=1}^a\left(1+R\left(\frac{E_i}{R^{A_i}}\right)^2\right)
   \prod_{j=1}^b\left(1+R\left(\frac{G_j}{R^{B_j}}\right)^2\right).
\end{equation}
Assume all displayed denominators are nonzero, $W\ne0$, and
\[
 T=\Tr_{F/\Q}(W^{-1})\ne0,\qquad \mathfrak b=R/T^2.
\]
The only property of a pole certificate needed in the following theorem
is the stated exclusion of negative odd valuations of $\mathfrak b$.

\begin{theorem}\label{thm:two-collections}
Assume $v_p(\mathfrak b)$ is not negative and odd.
\begin{enumerate}
\item If $v_p(z)=\nu>0$ is odd and $v_p(w)=0$, there is a place
$\mathfrak p\mid p$ such that
\[
              v_{\mathfrak p}(S_j)\ge n_j\nu+1
                       \qquad(1\le j\le b).
\]
At every place above $p$, each factor indexed by a $C_i$ in
\eqref{eq:product} is a unit.
\item If $v_p(w)=\nu>0$ is odd and $v_p(z)=0$, there is a place
$\mathfrak p\mid p$ such that
\[
              v_{\mathfrak p}(C_i)\ge m_i\nu+1
                       \qquad(1\le i\le a).
\]
At every place above $p$, each factor indexed by an $S_j$ is a unit.
\end{enumerate}
The conclusion holds without any bound on the valuations of the
collection whose factors are units.
\end{theorem}
\begin{proof}
We prove the first assertion; the second follows by interchanging
the two collections and $z,w$. Since $v_p(R)=\nu$ is odd,
the valuation of every factor in~\eqref{eq:product} is the minimum
of zero and the valuation of its second term. There is no cancellation
between terms of equal valuation: the second term has odd valuation.
Every factor therefore has nonpositive valuation, and so does $W$.

By~\eqref{eq:inactive}, $v_{\mathfrak p}(E_i)\ge K\nu$.
Consequently
\[
 v_{\mathfrak p}\left(R(E_i/R^{A_i})^2\right)
 \ge(1+2K-2A_i)\nu>0.
\]
Thus all factors in the first product are units.

Put $q=\max_{\mathfrak p\mid p}v_{\mathfrak p}(W)\le0$.
Local traces of integral elements are integral, so the decomposition
of the global trace into local traces gives $v_p(T)\ge-q$.
The integer
\[
                    v_p(\mathfrak b)=\nu-2v_p(T)
\]
is odd. By the assumed pole exclusion it is at least one, and hence
$q>-\nu/2$. Choose a place at which this maximum is attained.
Since the factor valuations are nonpositive and have sum $q$,
each one is greater than $-\nu/2$. For a factor indexed by $j$ this gives
\[
           2v_{\mathfrak p}(G_j)-(2B_j-1)\nu>-\nu/2,
\]
or
\[
              v_{\mathfrak p}(G_j)>(B_j-3/4)\nu>0.
\]
Here $w$ is a unit. Lemma~\ref{lem:normalization} shows that
$v_{\mathfrak p}(S_j)>0$ and
$v_{\mathfrak p}(G_j)=2v_{\mathfrak p}(S_j)$, unless $S_j=0$,
in which case the conclusion already holds. Therefore
\[
 v_{\mathfrak p}(S_j)>(n_j+1/8)\nu>n_j\nu.
\]
The valuation is an integer, proving the assertion.
\end{proof}

The trace argument does not require $p$ to split completely.
It uses only unramifiedness, which ensures that the odd/even valuation
comparison is valid at every place. Nor does it assume that the
local traces are units: cancellation in a trace only strengthens the
inequality used in the proof.

\section{A unit among two balanced contact functions}

Let $\Theta_0,\Theta_1$ be finite sets of nonzero rational numbers,
and let $c_0,c_1\in\Q^\times$ satisfy $c_i\notin\Theta_i$.
Assume
\[
       (\Theta_0\cup\{c_0\})\cap
       (\Theta_1\cup\{c_1\})=\varnothing.
\]
Put $d_i=|\Theta_i|$ and define the balanced rational functions
\[
                 R_i(t)=\frac{\prod_{\theta\in\Theta_i}(t-\theta)}
                                      {(t-c_i)^{d_i}}.
\]
The term balanced here means precisely that numerator and denominator
have the same degree; no additional property is included in the term.
If the intended application needs a square denominator, take $d_i$ even.
Let $S$ contain all prime divisors of the numerators and denominators
of the nonzero members of the two displayed sets and of their
nonzero pairwise differences.

\begin{theorem}\label{thm:unit}
Let $r,T,\lambda\in\Q$, with $r\ne0$, and set
\[
                  a=(1+rT^2)^2,\qquad t_0=\lambda,\quad t_1=a\lambda.
\]
Suppose $p\notin S$ and $v_p(r)$ is odd. If $R_i(t_i)$ are
defined and nonzero, then at least one is a $p$-adic unit:
\[
                       v_p(R_0(t_0))=0
                       \quad\hbox{or}\quad v_p(R_1(t_1))=0.
\]
The conclusion permits negative as well as positive odd $v_p(r)$.
\end{theorem}
\begin{proof}
If $T=0$, then $a=1$. Otherwise $v_p(rT^2)$ is odd and cannot
be zero. If it is positive, $a$ is a unit congruent to one modulo $p$.
If it is negative, $v_p(a)=2v_p(rT^2)<0$.
Thus in every case $a\ne0$ and either $a\equiv1\pmod p$ or
$v_p(a)<0$.

For either balanced contact function, an argument of negative valuation
has a numerator and denominator of the same valuation $d_i v_p(t)$,
so its value is a unit. An argument of positive valuation, including
zero, also gives a unit, since all the contacts and the pole are units.

It remains to consider $v_p(\lambda)=0$ and
$v_p(R_0(\lambda))\ne0$. The residue of $\lambda$ must then belong to
$\Theta_0\cup\{c_0\}$. If $v_p(a)<0$, the argument $a\lambda$ has
negative valuation, and $R_1(a\lambda)$ is a unit. Otherwise
$a\lambda\equiv\lambda\pmod p$. The two contact-and-pole sets have
disjoint reductions modulo $p$ by the choice of $S$, so this residue
does not belong to $\Theta_1\cup\{c_1\}$. Every numerator and denominator
factor of $R_1(a\lambda)$ is then a unit.
\end{proof}

\section{Scope for Diophantine constructions}

Theorem~\ref{thm:unit} supplies one contact function of valuation exactly
zero at a prime of odd valuation of $r$. Thus
Theorem~\ref{thm:two-collections} can extract scalar local conditions
from at least one product $rR_i(t_i)$, while its other collection of
local conditions is suppressed independently of any pole depths.
At a comparison prime where $r$ is a unit and $R_i(t_i)$ has positive
odd valuation, the same theorem instead extracts the comparison
conditions. Every transformation is rational and introduces no new
unrestricted rational parameter; its denominators must be retained
as nonzero guards in an algebraic presentation.

These local facts do not suffice for a global defining-polynomial
theorem. In particular, they give no extraction statement at a prime
where both coefficients have nonzero valuation of the same parity.
Deleting such primes from a contact radical changes the global height
estimate. Their total contribution is not bounded by a constant when
$r$ depends on the rational input. A construction which combines
scalar and comparison pole inputs must also prove its global height
inequalities and its simultaneous rational realization. No reduction
in the number of defining-polynomial variables is asserted here.
\end{document}
